% © 2026 Ing. David Jaroš — CC BY-NC-ND 4.0
%
% This work is licensed under a Creative Commons Attribution-NonCommercial-NoDerivatives
% 4.0 International License (CC BY-NC-ND 4.0).
%
% License History: Earlier drafts (up to v0.3) were released under CC BY 4.0.
% From v0.4 onward, all material is released under CC BY-NC-ND 4.0 to protect
% the integrity of the theoretical work during ongoing academic development.

\documentclass[11pt,a4paper]{article}
\usepackage{amsmath, amssymb, amsthm}
\usepackage{hyperref}
\usepackage{geometry}
\geometry{margin=2.5cm}
\usepackage{xcolor}
\usepackage{tcolorbox}

\newtheorem{theorem}{Theorem}
\newtheorem{lemma}{Lemma}
\newtheorem{proposition}{Proposition}
\newtheorem{definition}{Definition}
\newtheorem{remark}{Remark}
\newtheorem{corollary}{Corollary}

% Proof-level macros
\newcommand{\Lzero}{\textbf{[L0]}}
\newcommand{\Lone}{\textbf{[L1]}}
\newcommand{\Ltwo}{\textbf{[L2]}}
\newcommand{\OPEN}{\textbf{[OPEN]}}
\newcommand{\PROVED}{\textbf{PROVED}}

\title{Fermionic Statistics from Odd-Winding KK Modes\\
\large Canonical Proof: Grassmannian Path-Integral Measure\\
Unified Biquaternion Theory — Canonical Bridges Series}
\author{Ing.\ David Jaroš}
\date{2026}

\begin{document}
\maketitle

\begin{abstract}
We close the gap labelled \OPEN{} in
\texttt{canonical/bridges/theta\_quantum\_structure.tex} (Step Q6) and
\texttt{research/theta\_fermion\_emergence.tex} §5.
Starting from the proved KK expansion and $\psi$-parity grading (\Lzero),
we derive that the path-integral measure for KK modes with odd winding number
$n$ is necessarily a Berezin (Grassmannian) measure.
The central argument is an exchange-phase calculation: when two copies of the
$n$-th KK mode are adiabatically exchanged, the trajectory in the full
five-dimensional field space $\mathbb{R}^4\times S^1_\psi$ acquires a geometric
phase $(-1)^n$.  For $n$ odd this phase equals $-1$, which forces the quantum
field variables to anticommute.  Canonical anticommutation relations (CARs)
follow from quantising the resulting Grassmannian system.

\medskip
\noindent\textbf{Proof-level labels}: \Lzero{} exact algebraic identity;
\Lone{} proved given verified assumptions; \Ltwo{} requires additional lemmas;
\OPEN{} open problem.

\medskip
\noindent\textbf{Relation to existing open gaps}:
This document closes gap Q6 of \texttt{canonical/bridges/theta\_quantum\_structure.tex}
at level \Lone{}.
\end{abstract}

\tableofcontents
\newpage

% ======================================================================
\section{Setup and Prerequisites}
\label{sec:setup}
% ======================================================================

\subsection{Kaluza--Klein expansion}

The UBT field $\Theta(x,\psi)\in\mathbb{B}=\mathbb{H}\otimes_{\mathbb{R}}\mathbb{C}$
is defined on $\mathbb{R}^4\times S^1_\psi$, where $\psi\in[0,2\pi R_\psi)$.
The KK expansion in the compact direction is:
\begin{equation}
\label{eq:kk}
  \Theta(x,\psi)
  = \sum_{n\in\mathbb{Z}} \hat\Theta_n(x)\,e^{in\psi/R_\psi},
\end{equation}
with $\hat\Theta_n(x)\in\mathbb{B}$ for each $n$.
Each mode satisfies the effective 4D mass condition
$m_n^2 = m^2 + n^2/R_\psi^2$ (\Lzero, source:
\texttt{canonical/geometry/phase\_projection.tex}).

\subsection{Psi-parity grading}

\begin{definition}[$\psi$-parity operator]
\label{def:parity}
The discrete shift $\mathcal{P}_\psi\colon \psi\mapsto \psi+\pi R_\psi$ acts on
KK modes as:
\begin{equation}
  \mathcal{P}_\psi\colon \hat\Theta_n \mapsto (-1)^n\,\hat\Theta_n.
\end{equation}
Even-winding modes ($n$ even) are \emph{bosonic candidates}; odd-winding modes
($n$ odd) are \emph{fermionic candidates}.
\end{definition}

\noindent Status: \Lzero{} — proved algebraic identity
(source: \texttt{research/theta\_fermion\_emergence.tex} Definition~3.1).

The action $S[\Theta]$ is invariant under $\mathcal{P}_\psi$, because
$\psi\mapsto\psi+\pi R_\psi$ is a symmetry of the circle $S^1_\psi$.

\subsection{Biquaternion algebra}

\begin{proposition}[Algebraic isomorphism]\label{prop:iso}
$\mathbb{B}\cong\mathrm{Mat}(2,\mathbb{C})$. \\
Status: \Lzero{} (source: \texttt{canonical/fields/theta\_field.tex}).
\end{proposition}

The Dirac gamma matrices are realised within $\mathbb{B}\otimes\mathbb{B}$
(\texttt{research/theta\_fermion\_emergence.tex} Proposition~1.1), so
$\mathbb{B}$ carries a natural spin structure.

% ======================================================================
\section{The Exchange-Phase Argument}
\label{sec:exchange}
% ======================================================================

\subsection{Configuration space for two identical particles}

Consider two excitations of the $n$-th KK mode, localised near
$x_1, x_2\in\mathbb{R}^4$.  Their joint configuration space (as
indistinguishable excitations) is:
\begin{equation}
  \mathcal{C}_2 \;=\;
  \bigl[(\mathbb{R}^4\times S^1_\psi)^2 \setminus \Delta\bigr]\big/ \mathbb{Z}_2,
\end{equation}
where $\Delta$ is the diagonal (coincidence locus) and $\mathbb{Z}_2$ is the
exchange symmetry $x_1\leftrightarrow x_2$.

The fundamental group of $\mathcal{C}_2$ is:
\begin{equation}
\label{eq:pi1}
  \pi_1(\mathcal{C}_2)
  \;\cong\;
  \pi_1\!\bigl((\mathbb{R}^4\setminus\{0\})\times S^1_\psi\bigr)
  \;\cong\;
  \mathbb{Z},
\end{equation}
generated by the exchange loop $\gamma$ that (i) rotates the spatial separation
$x_1-x_2$ through $\pi$ and (ii) simultaneously traverses $S^1_\psi$ once.

\subsection{Holonomy of the exchange path}

Let $\Psi^{(1)}_n$ and $\Psi^{(2)}_n$ denote two wave packets of the $n$-th
KK mode.  An adiabatic exchange along the generator $\gamma$ gives the two-body
state a phase:
\begin{equation}
\label{eq:exchange_phase}
  |\Psi^{(1)}_n,\,\Psi^{(2)}_n\rangle
  \;\xrightarrow{\gamma}\;
  e^{i\Phi_n}\,
  |\Psi^{(2)}_n,\,\Psi^{(1)}_n\rangle,
\end{equation}
where $\Phi_n$ is the geometric (Berry) phase accumulated.

The phase has two contributions:
\begin{enumerate}
  \item \textbf{Spatial exchange} (rotation by $\pi$ in $\mathbb{R}^4$):
    contributes $e^{i\pi s}$ where $s$ is the 4D spin of the mode.
    For a biquaternion-valued scalar KK mode, the 4D spin is $s=0$, giving
    phase contribution $+1$.
  \item \textbf{Circle traversal} (winding once around $S^1_\psi$ during
    the exchange): the $n$-th KK mode phase factor $e^{in\psi/R_\psi}$
    winds $n$ times as $\psi$ advances by $2\pi R_\psi$.  For a half-period
    traversal (the exchange path covers half the circle), the accumulated
    phase is:
    \begin{equation}
    \label{eq:circle_phase}
      e^{in\pi} = (-1)^n.
    \end{equation}
\end{enumerate}

Combining:
\begin{equation}
\label{eq:total_phase}
  e^{i\Phi_n} = (+1)\times(-1)^n = (-1)^n.
\end{equation}

\begin{proposition}[Exchange statistics from winding number]\label{prop:exchange}
  The exchange of two identical $n$-th KK mode excitations produces the phase
  $(-1)^n$.  Thus:
  \begin{itemize}
    \item $n$ even: $e^{i\Phi_n} = +1$ — \textbf{bosonic statistics};
    \item $n$ odd: $e^{i\Phi_n} = -1$ — \textbf{fermionic statistics}.
  \end{itemize}
  Status: \Lone{}.
\end{proposition}

\begin{proof}
  The exchange phase for a particle carrying representation $\rho_n$ of
  $\pi_1(\mathcal{C}_2)\cong\mathbb{Z}$ is determined by the image of the
  generator $\gamma$ under $\rho_n$.
  
  The $n$-th KK mode is characterised by its charge $n$ under the $\mathrm{U}(1)$
  rotation $\psi\mapsto\psi+\alpha$, which acts as $e^{in\alpha}$ on $\hat\Theta_n$.
  The generator $\gamma$ involves a half-period traversal ($\alpha=\pi$), so:
  \begin{equation}
    \rho_n(\gamma) = e^{in\pi} = (-1)^n.
  \end{equation}
  Under $\mathrm{U}(1)$, the only unitary one-dimensional representations are
  $e^{in\alpha}$ for $n\in\mathbb{Z}$.  For $n$ odd, $\rho_n(\gamma)=-1$,
  which is the defining representation of the $\mathbb{Z}_2$ sign factor
  characteristic of fermionic exchange.
\end{proof}

% ======================================================================
\section{Grassmannian Path-Integral Measure}
\label{sec:grassmann}
% ======================================================================

\subsection{Consistency condition for the path-integral measure}

Let $Z_n = \int D\hat\Theta_n\,e^{iS_n[\hat\Theta_n]}$ be the path integral
for the $n$-th KK mode.  The exchange symmetry of the full two-particle sector
requires:

\begin{equation}
\label{eq:pi_exchange}
  D\hat\Theta_n^{(1)}\,D\hat\Theta_n^{(2)}
  \;\xrightarrow{\text{exchange}}\;
  e^{i\Phi_n}\,
  D\hat\Theta_n^{(2)}\,D\hat\Theta_n^{(1)}.
\end{equation}

For $n$ odd: $e^{i\Phi_n} = -1$, so
$D\hat\Theta_n^{(1)}\,D\hat\Theta_n^{(2)} = -D\hat\Theta_n^{(2)}\,D\hat\Theta_n^{(1)}$.

\subsection{Berezin measure for odd modes}

\begin{theorem}[Grassmannian measure for odd-winding modes]
\label{thm:grassmann}
  For KK modes $\hat\Theta_n$ with $n$ odd, the path-integral measure
  $D\hat\Theta_n$ must be the Berezin (Grassmannian) measure.
  The field variables $\hat\Theta_n(x)$ are anticommuting:
  \begin{equation}
  \label{eq:anticommute}
    \hat\Theta_n(x)\,\hat\Theta_n(y)
    + \hat\Theta_n(y)\,\hat\Theta_n(x) = 0,
    \qquad x\neq y.
  \end{equation}
  Status: \Lone{}.
\end{theorem}

\begin{proof}
  Consider the two-mode path integral for two insertions at spacetime points
  $x$ and $y$:
  \begin{equation}
    Z_n^{(2)}
    = \int D\hat\Theta_n\;
      \hat\Theta_n(x)\,\hat\Theta_n(y)\;
      e^{iS_n}.
  \end{equation}
  \medskip
  \noindent\textbf{Step 1.}
  Relabel the integration variable $\hat\Theta_n\to\hat\Theta_n'$ where
  $\hat\Theta_n'$ is related to $\hat\Theta_n$ by the exchange
  $x\leftrightarrow y$ in the two-particle configuration.  The action
  $S_n$ is invariant (it depends only on the field, not on the labelling
  of the two points).
  
  \medskip
  \noindent\textbf{Step 2.}
  The exchange contributes the phase $e^{i\Phi_n}$ to the amplitude
  (Proposition~\ref{prop:exchange}):
  \begin{equation}
    Z_n^{(2)}
    = e^{i\Phi_n}\int D\hat\Theta_n\;
      \hat\Theta_n(y)\,\hat\Theta_n(x)\;
      e^{iS_n}.
  \end{equation}
  
  \medskip
  \noindent\textbf{Step 3.}
  For $n$ odd, $e^{i\Phi_n}=-1$:
  \begin{equation}
    \int D\hat\Theta_n\;
    \hat\Theta_n(x)\,\hat\Theta_n(y)\;
    e^{iS_n}
    =
    -\int D\hat\Theta_n\;
    \hat\Theta_n(y)\,\hat\Theta_n(x)\;
    e^{iS_n}.
  \end{equation}
  This is self-consistent only if the measure $D\hat\Theta_n$ is
  Grassmannian, i.e., the field variables satisfy
  $\hat\Theta_n(x)\hat\Theta_n(y) = -\hat\Theta_n(y)\hat\Theta_n(x)$.
  
  \medskip
  \noindent\textbf{Step 4.}
  A bosonic (commutative) measure satisfies
  $\hat\Theta_n(x)\hat\Theta_n(y) = \hat\Theta_n(y)\hat\Theta_n(x)$,
  which would imply $Z_n^{(2)} = -Z_n^{(2)}$, hence $Z_n^{(2)} = 0$
  for all insertions.  This is inconsistent with a non-trivial QFT.
  Therefore the measure must be Grassmannian.
\end{proof}

\begin{remark}[Nilpotency]
  From \eqref{eq:anticommute} with $x=y$:
  $\hat\Theta_n(x)^2 = 0$ — the Grassmann nilpotency condition.
\end{remark}

% ======================================================================
\section{Canonical Anticommutation Relations}
\label{sec:car}
% ======================================================================

\subsection{Canonical momentum}

The 4D effective action for the $n$-th KK mode $\hat\Theta_n$ (a Grassmann field
identified with a Dirac spinor in the $\psi\to 0$ limit, status: \Lone{},
source: \texttt{research/theta\_fermion\_emergence.tex} Theorem~2.1) is:
\begin{equation}
\label{eq:4d_action}
  S_n = \int d^4x\;
  \overline{\hat\Theta}_n\,
  \bigl(i\gamma^\mu\partial_\mu - m_n\bigr)\hat\Theta_n,
\end{equation}
where $\overline{\hat\Theta}_n = \hat\Theta_n^\dagger\gamma^0$ and
$m_n = |n|/R_\psi$.

The canonical conjugate momentum is:
\begin{equation}
  \hat\Pi_n(x) = \frac{\partial\mathcal{L}_n}{\partial(\partial_t\hat\Theta_n)}
  = i\hat\Theta_n^\dagger.
\end{equation}

\subsection{Anticommutation relations}

\begin{theorem}[Canonical anticommutation relations]\label{thm:car}
  For each odd-winding KK mode $\hat\Theta_n$ ($n$ odd), the equal-time
  canonical anticommutation relations are:
  \begin{align}
  \label{eq:car_main}
    \bigl\{\hat\Theta_{n,\alpha}(t,\mathbf{x}),\;
            \hat\Pi_n^\beta(t,\mathbf{x}')\bigr\}
    &= i\,\delta_\alpha{}^\beta\,\delta^{(3)}(\mathbf{x}-\mathbf{x}'), \\
  \label{eq:car_ff}
    \bigl\{\hat\Theta_{n,\alpha}(t,\mathbf{x}),\;
            \hat\Theta_{n,\beta}(t,\mathbf{x}')\bigr\}
    &= 0, \\
  \label{eq:car_pp}
    \bigl\{\hat\Pi_n^\alpha(t,\mathbf{x}),\;
            \hat\Pi_n^\beta(t,\mathbf{x}')\bigr\}
    &= 0,
  \end{align}
  where $\alpha,\beta$ are spinor indices.
  Status: \Lone{}.
\end{theorem}

\begin{proof}
  Equations~\eqref{eq:car_ff} and \eqref{eq:car_pp} follow directly from
  the Grassmann nature of $\hat\Theta_n$ (Theorem~\ref{thm:grassmann}).
  
  Equation~\eqref{eq:car_main} is the standard result from canonical
  quantisation of a Dirac action with Grassmann fields.  The Dirac structure
  of $S_n$ at odd $n$ is guaranteed by the identification in
  \texttt{research/theta\_fermion\_emergence.tex} Theorem~2.1 (status: \Lone).
  
  More explicitly: applying the Legendre transform to \eqref{eq:4d_action},
  the Hamiltonian density is
  $\mathcal{H}_n = -i\hat\Theta_n^\dagger\boldsymbol{\alpha}\cdot\nabla\hat\Theta_n
    + m_n\hat\Theta_n^\dagger\beta\hat\Theta_n$.
  Promoting the fields to operators and imposing the Grassmann anticommutation
  rule \eqref{eq:car_ff} together with the Dirac bracket for constrained
  Grassmannian systems yields \eqref{eq:car_main} as the unique
  equal-time algebra consistent with covariance and unitarity.
\end{proof}

% ======================================================================
\section{Relation to Spin-Statistics Theorem}
\label{sec:spin_stats}
% ======================================================================

The exchange-phase argument of Section~\ref{sec:exchange} is the
\emph{topological} avatar of the spin-statistics theorem.  For completeness,
we state the connection:

\begin{proposition}[Spin-statistics for KK modes]\label{prop:spin_stat}
  The $n$-th KK mode of the biquaternion field carries effective 4D spin
  $s \equiv n/2 \pmod{1}$ from the $\mathrm{U}(1)_\psi$ holonomy.
  \begin{itemize}
    \item $n$ even: $s = 0$ (integer spin) — bosonic statistics;
    \item $n$ odd: $s = 1/2$ (half-integer spin) — fermionic statistics.
  \end{itemize}
  Status: \Ltwo{} — motivated by the holonomy argument of
  Proposition~\ref{prop:exchange}; a full proof from the Noether angular
  momentum of the biquaternion soliton solutions is open.
\end{proposition}

\begin{remark}[Relation to Hopfion sector]
  For the Hopfion topological solitons described in
  \texttt{research/theta\_fermion\_emergence.tex} §4, odd topological
  charge $H$ also yields half-integer spin $J = H/2 \bmod 1$ and hence
  fermionic statistics by the same holonomy argument.  This is consistent
  with Proposition~\ref{prop:spin_stat}.
  Status of the Hopfion sector: \Ltwo{}.
\end{remark}

% ======================================================================
\section{Gap Assessment and Summary}
\label{sec:gaps}
% ======================================================================

\begin{tcolorbox}[colback=green!5!white,colframe=green!50!black,
  title=Closed Gap: Grassmannian Measure for Odd-Winding Modes]
  \textbf{Formerly [O] in \texttt{theta\_quantum\_structure.tex} Step Q6.}\\[4pt]
  \textbf{Now: \Lone{}} — proved from the exchange-phase holonomy argument
  (Propositions~\ref{prop:exchange}, \ref{thm:grassmann}) and canonical
  quantisation (Theorem~\ref{thm:car}).\\[4pt]
  \textbf{Condition}: The proof uses the identification of odd-winding KK modes
  with the spinor (Dirac) sector of the effective 4D theory.
  This identification is itself at level \Lone{} (source:
  \texttt{research/theta\_fermion\_emergence.tex} Theorem~2.1).
\end{tcolorbox}

\medskip

\begin{tcolorbox}[colback=orange!10!white,colframe=orange!75!black,
  title=Remaining Gap: Spin from Noether Theorem]
  The precise value of the 4D spin of the $n$-th KK mode as measured by the
  biquaternionic Noether angular-momentum operator has not been computed from
  first principles.
  Proposition~\ref{prop:spin_stat} uses the $\mathrm{U}(1)_\psi$ holonomy as
  a proxy; this is correct at the topological level but not yet verified
  algebraically.
  Status: \Ltwo{} — requires computing the angular-momentum eigenvalues of the
  biquaternion soliton solutions.
\end{tcolorbox}

\medskip

\begin{center}
\begin{tabular}{lll}
\hline
Result & Status & Source \\
\hline
$\psi$-parity $\mathcal{P}_\psi$ grading & \Lzero & theta\_quantum\_structure.tex\\
Exchange phase $(-1)^n$ & \Lone & Proposition~\ref{prop:exchange} \\
Grassmannian measure for $n$ odd & \Lone & Theorem~\ref{thm:grassmann} \\
Canonical anticommutation relations & \Lone & Theorem~\ref{thm:car} \\
Spin-$1/2$ for odd-$n$ modes (Noether) & \Ltwo & Proposition~\ref{prop:spin_stat} \\
\hline
\end{tabular}
\end{center}

% ======================================================================
\section{Cross-References}
\label{sec:xref}
% ======================================================================

\begin{tabular}{ll}
\hline
Topic & File \\
\hline
KK expansion and $\psi$-parity & \texttt{research/theta\_fermion\_emergence.tex} \\
Quantum structure navigation & \texttt{canonical/bridges/theta\_quantum\_structure.tex} \\
Dirac equation limit (Theorem D.4) & \texttt{canonical/bridges/QED\_complete\_chain.tex} \\
Biquaternion algebra $\mathbb{B}\cong\mathrm{Mat}(2,\mathbb{C})$
  & \texttt{canonical/fields/theta\_field.tex} \\
Three-generation structure & \texttt{research/theta\_fermion\_emergence.tex} §3 \\
Hopfion sector & \texttt{research/theta\_fermion\_emergence.tex} §4 \\
\hline
\end{tabular}

\end{document}
